\subsection{Desafios Técnicos}
Durante o desenvolvimento deste trabalho, um dos principais desafios encontrados foi o tratamento do desbalanceamento presente no conjunto de dados, uma vez que essa característica pode influenciar diretamente o treinamento e a avaliação dos classificadores. Para mitigar esse problema, foi empregada a técnica SMOTE, permitindo analisar os impactos do balanceamento artificial das classes no desempenho dos modelos.

Outro desafio relevante esteve relacionado à elevada dimensionalidade dos dados e à identificação dos atributos mais representativos para a tarefa de classificação. A aplicação de técnicas de redução de dimensionalidade, bem como técnicas de visualização de possíveis correlações entre atributos e métodos de agrupamento não supervisionados, permitiu avaliar a influência da simplificação do espaço de atributos sobre a capacidade de generalização dos algoritmos. Além disso, foi necessário realizar o ajuste de hiperparâmetros por meio de busca sistemática, garantindo que cada modelo fosse avaliado em sua configuração mais adequada.

\subsection{Conclusões Finais}
O presente trabalho teve como objetivo avaliar o desempenho de diferentes algoritmos de aprendizado de máquina na tarefa de predição da evasão acadêmica, comparando os efeitos do balanceamento de classes por meio do SMOTE e da redução de dimensionalidade sobre os resultados obtidos. Os experimentos realizados demonstraram que os modelos \textit{Support Vector Machine} e \textit{Multi-Layer Perceptron} apresentaram os melhores desempenhos gerais, alcançando acurácia de até 0,76 e $F_1$-Score próximo de 0,70.

Os resultados indicaram que a redução de dimensionalidade foi a estratégia de pré-processamento mais eficaz para o conjunto de dados estudado, proporcionando melhorias ou manutenção do desempenho para a maioria dos classificadores avaliados. Em contrapartida, a aplicação do SMOTE não resultou em ganhos significativos e, em diversos cenários, ocasionou redução das métricas de desempenho. Esses resultados sugerem que a eliminação de atributos redundantes ou pouco relevantes exerceu maior influência na qualidade da classificação do que o balanceamento artificial das classes.

Ademais, ao confrontar os achados deste estudo com a literatura correlacionada, observa-se que os dois classificadores de melhor desempenho obtiveram índices de acurácia superiores aos reportados na referência~[4]. Embora uma comparação direta estrita seja limitada pelas inevitáveis divergências nos cenários experimentais e nos modelos desenvolvidos subjacentes, esses resultados sugerem a robustez e a eficácia das estratégias de modelagem e pré-processamento adotadas neste trabalho, sinalizando o potencial dos modelos propostos para o contexto avaliado.

Dessa forma, conclui-se que a seleção e transformação adequadas dos atributos desempenham papel fundamental na construção de modelos preditivos para evasão acadêmica. Como trabalhos futuros, sugere-se a investigação de outras técnicas de balanceamento de dados, métodos mais avançados de seleção de atributos e algoritmos aplicados, bem como a avaliação de arquiteturas neurais mais complexas e estratégias de otimização adicionais.
